% =====================================================================
% Modelo ABNT de trabalho academico (TCC, monografia, dissertacao, tese)
% EguaTeX (www.eguatex.com.br). Licenca CC0 1.0: dominio publico.
%
% Segue a ABNT NBR 14724:2024 e as normas que ela chama: NBR 6023:2025
% (referencias), NBR 6024:2012 (secoes), NBR 6027:2012 (sumario),
% NBR 6028:2021 (resumo) e NBR 10520:2023 (citacoes).
%
% Como usar:
%   1. Preencha os dados do trabalho logo abaixo.
%   2. Troque cada texto de exemplo pelo seu conteudo.
%   3. Apague o que for opcional e voce nao quiser usar.
%   4. Coloque as obras que voce citar no arquivo referencias.bib.
%   5. Compile com pdfLaTeX e BibTeX (o EguaTeX faz isso sozinho).
% A formatacao fica toda no arquivo abnt-eguatex.sty. Nao precisa mexer nele.
%
% Antes de entregar, confira o manual da sua instituicao: ela pode pedir
% detalhes que a ABNT deixa livres (fonte, lombada, modelo de capa).
% =====================================================================
\documentclass[12pt,a4paper,oneside]{article}
\usepackage[academico]{abnt-eguatex}

% ---------------------------------------------------------------------
% DADOS DO TRABALHO
% ---------------------------------------------------------------------
\instituicao{Nome da instituição}   % opcional na capa; deixe {} vazio para tirar
\autor{Nome do autor}
\titulo{Título do trabalho}
\subtitulo{subtítulo, se houver}    % deixe {} vazio se nao houver
\local{Cidade}
\ano{2026}                          % ano da entrega
% Natureza: tipo do trabalho, objetivo, instituicao e area de concentracao.
% Para dissertacao ou tese, troque o tipo e o grau pretendido.
\natureza{Trabalho de conclusão de curso apresentado ao Curso de Nome do Curso
  da Nome da Instituição, como requisito parcial para obtenção do grau de
  Bacharel em Nome do Curso.\par\vspace{\baselineskip}
  Área de concentração: nome da área.}
\orientador{Orientador(a): Prof.(ª) Dr.(ª) Nome do orientador}
\coorientador{Coorientador(a): Prof.(ª) Dr.(ª) Nome do coorientador} % apague se nao houver

\begin{document}

% =====================================================================
% PARTE EXTERNA
% =====================================================================
\imprimircapa          % obrigatoria (NBR 14724:2024, 4.1.1)

% =====================================================================
% ELEMENTOS PRE-TEXTUAIS (contados, mas sem numero na pagina)
% =====================================================================
\imprimirfolhaderosto  % obrigatoria (4.2.1.1)
\imprimirficha         % verso da folha de rosto; a biblioteca faz a ficha

% Errata (opcional, 4.2.1.2): so existe depois do trabalho impresso.
% Tire os sinais de % das linhas abaixo para usar.
% \titulopretextual{Errata}
% \noindent SOBRENOME, Prenome. \textbf{Título do trabalho}: subtítulo. 2026. Trabalho de Conclusão de Curso (Bacharelado em Nome do Curso) -- Nome da Instituição, Cidade, 2026.
% \par\vspace{\baselineskip}
% \begin{tabular}{llll}\toprule Folha & Linha & Onde se lê & Leia-se \\ \midrule 16 & 10 & texo & texto \\ \bottomrule\end{tabular}

% Folha de aprovacao (obrigatoria, 4.2.1.3). Data e assinaturas so depois
% da defesa.
\begin{folhadeaprovacao}
  \membrodabanca{Prof.(ª) Dr.(ª) Nome do orientador (orientador)}{Nome da Instituição}
  \membrodabanca{Prof.(ª) Dr.(ª) Nome do avaliador}{Nome da Instituição}
  \membrodabanca{Prof.(ª) Dr.(ª) Nome do avaliador}{Nome da Instituição}
\end{folhadeaprovacao}

% Dedicatoria (opcional, 4.2.1.4)
\begin{dedicatoria}
  Aqui vai a dedicatória, se você quiser. Ela não tem título.
\end{dedicatoria}

% Agradecimentos (opcional, 4.2.1.5)
\begin{agradecimentos}
  Aqui vão os agradecimentos a quem contribuiu de forma relevante para o trabalho.
\end{agradecimentos}

% Epigrafe (opcional, 4.2.1.6): uma citacao ligada ao tema, com a autoria.
\begin{epigrafe}
  Aqui vai a epígrafe, uma citação relacionada ao tema do trabalho.
  \autoriadaepigrafe{Nome do autor da citação}
\end{epigrafe}

% Resumo na lingua do texto (obrigatorio, 4.2.1.7; NBR 6028:2021)
% Paragrafo unico, verbo na terceira pessoa, de 150 a 500 palavras.
\begin{resumo}
Aqui vai o resumo do trabalho, escrito em um único parágrafo, com frases curtas e o verbo na terceira pessoa. Ele apresenta o objetivo, o método, os resultados e as conclusões, de modo que o leitor entenda o trabalho sem precisar ler o texto inteiro. Em trabalhos acadêmicos, o resumo tem de 150 a 500 palavras. Evite símbolos, fórmulas e abreviações que não sejam de uso comum.
\palavraschave{palavra-chave; segunda palavra-chave; terceira palavra-chave.}
\end{resumo}

% Resumo em lingua estrangeira (obrigatorio, 4.2.1.8)
\begin{otherlanguage}{english}
\begin{resumo}[Abstract]
Here goes the abstract, the translation of the summary into a language of international circulation, such as English. It follows the same rules as the summary in Portuguese: a single paragraph, concise sentences and the same content.
\palavraschave[Keywords]{keyword; second keyword; third keyword.}
\end{resumo}
\end{otherlanguage}

% Listas (opcionais, 4.2.1.9 a 4.2.1.12). Apague as que nao usar.
\listadefiguras
\listadequadros
\listadetabelas

\begin{listadesiglas}
  \sigla{ABNT}{Associação Brasileira de Normas Técnicas}
  \sigla{IBGE}{Instituto Brasileiro de Geografia e Estatística}
  \sigla{NBR}{Norma Brasileira}
\end{listadesiglas}

\begin{listadesiglas}[Lista de símbolos]
  \sigla{$v$}{Velocidade}
  \sigla{$\Delta t$}{Intervalo de tempo}
\end{listadesiglas}

% Sumario (obrigatorio, 4.2.1.13; NBR 6027:2012). E montado sozinho.
\sumario

% =====================================================================
% ELEMENTOS TEXTUAIS (a partir daqui as paginas tem numero)
% O nome das secoes fica a seu criterio. O trabalho e dividido em secoes,
% nao em capitulos (4.2.2).
% =====================================================================
\textual

\section{Introdução}

Aqui começa o texto. A introdução apresenta o tema, o objetivo do trabalho e as razões que levaram à pesquisa. O texto usa letra tamanho 12 e espaçamento de 1,5 entre as linhas.

Quando uma sigla aparece pela primeira vez, escreva o nome completo e a sigla entre parênteses, como em Associação Brasileira de Normas Técnicas (ABNT). Depois, use só a sigla.

\section{Desenvolvimento}

O desenvolvimento detalha a pesquisa. Ele pode ser dividido em quantas seções forem necessárias, até o quinto nível, como nos exemplos abaixo. Toda seção precisa ter algum texto.

\subsection{Seção secundária}

Texto da seção secundária. As citações seguem o sistema autor-data: o sobrenome vem em letras maiúsculas e minúsculas, seguido do ano e da página, como em \cite[p.~15]{livro2024}. Quando o nome do autor faz parte da frase, só o ano e a página ficam entre parênteses, como em \citeonline[p.~15]{livro2024}. No código, a citação é feita com \texttt{\textbackslash cite\{chave\}}, e a obra entra sozinha na lista de referências.

Uma citação direta de até três linhas fica no próprio parágrafo, entre aspas duplas: “aqui vai o trecho copiado da obra” \cite[p.~32]{livro2024}. A citação direta com mais de três linhas fica em destaque, como no exemplo a seguir.

\begin{citacao}
Aqui vai uma citação direta com mais de três linhas. Ela tem recuo de 4 cm da margem esquerda, letra menor que a do texto, espaçamento simples e não leva aspas. A indicação da fonte vem no final, entre parênteses, com o sobrenome do autor, o ano e a página de onde o trecho foi tirado \cite[p.~47]{capitulo2023}.
\end{citacao}

Uma citação indireta resume a ideia do autor com as suas palavras e não precisa de aspas \cite{artigo2025}. Obras de vários tipos entram na lista do mesmo jeito, como um trabalho de conclusão de curso \cite{tcc2022} ou uma página da internet \cite{pagina2025}. Notas de rodapé servem para explicações que não cabem no texto\footnote{Aqui vai a nota de rodapé. Ela tem letra menor, espaçamento simples e fica separada do texto por um traço de 5 cm. Quando a nota passa de uma linha, a segunda linha começa embaixo da primeira letra da nota.}.

\subsubsection{Seção terciária}

Assuntos sem título próprio dentro de uma seção podem ser divididos em alíneas, e o texto que vem antes delas termina em dois-pontos:
\begin{alineas}
  \item a alínea começa com letra minúscula e termina em ponto e vírgula;
  \item a alínea que tem subalíneas termina em dois-pontos:
  \begin{subalineas}
    \item a subalínea começa com travessão e termina em ponto e vírgula;
    \item a última subalínea termina em ponto e vírgula se ainda houver alínea depois;
  \end{subalineas}
  \item a última alínea termina em ponto.
\end{alineas}

\paragraph{Seção quaternária}

Texto da seção quaternária.

\subparagraph{Seção quinária}

Texto da seção quinária, o último nível permitido.

\subsection{Ilustrações, tabelas e equações}

Toda ilustração é citada no texto, como a Figura~\ref{fig:exemplo}, e fica o mais perto possível do trecho em que é mencionada. Em cima vêm o nome, o número, um travessão e o título. Embaixo vem a fonte, sempre, mesmo quando a ilustração é sua.

\begin{figure}[htbp]
  \begin{ilustracao}{0.75\textwidth}
    \caption{Título da figura}\label{fig:exemplo}
    \fbox{\parbox[c][5cm][c]{\dimexpr\linewidth-2\fboxsep-2\fboxrule}{\centering Coloque a imagem aqui com\\\texttt{\textbackslash includegraphics[width=\textbackslash linewidth]\{arquivo\}}}}
    \fonte{elaborada pelo autor (2026).}
    \nota{aqui vai uma nota sobre a figura, se for preciso.}
  \end{ilustracao}
\end{figure}

Quadros organizam informações em texto, com linhas fechando as laterais, como o Quadro~\ref{qua:exemplo}.

\begin{quadro}[htbp]
  \begin{ilustracao}{0.75\textwidth}
    \caption{Título do quadro}\label{qua:exemplo}
    \setstretch{1}
    \begin{tabular}{|p{0.3\linewidth}|p{0.6\linewidth}|}
      \hline
      \textbf{Coluna 1} & \textbf{Coluna 2} \\ \hline
      Informação & Descrição em texto \\ \hline
      Informação & Descrição em texto \\ \hline
    \end{tabular}
    \fonte{\citeonline[p.~20]{livro2024}.}
  \end{ilustracao}
\end{quadro}

Tabelas apresentam dados numéricos e seguem as normas de apresentação tabular do IBGE: traços horizontais em cima, embaixo e abaixo do cabeçalho, sem traços nas laterais, como a Tabela~\ref{tab:exemplo}.

\begin{table}[htbp]
  \begin{ilustracao}{0.75\textwidth}
    \caption{Título da tabela, com o que, onde e quando}\label{tab:exemplo}
    \setstretch{1}\centering
    \begin{tabular*}{\linewidth}{@{\extracolsep{\fill}}lrr}
      \toprule
      Categoria & Quantidade & Percentual (\%) \\
      \midrule
      Categoria A & 120 & 60,0 \\
      Categoria B & 80 & 40,0 \\
      \midrule
      Total & 200 & 100,0 \\
      \bottomrule
    \end{tabular*}
    \fonte{elaborada pelo autor (2026).}
  \end{ilustracao}
\end{table}

Equações ficam em destaque e, se preciso, numeradas à direita, entre parênteses, como a Equação~\eqref{eq:exemplo}.
\begin{equation}
  v = \frac{\Delta s}{\Delta t}
  \label{eq:exemplo}
\end{equation}

\section{Conclusão}

A conclusão retoma os objetivos e apresenta os resultados alcançados.

% =====================================================================
% ELEMENTOS POS-TEXTUAIS
% =====================================================================

% Referencias (obrigatorias; NBR 6023:2025). A lista e montada sozinha
% a partir do arquivo referencias.bib: entram so as obras citadas no texto,
% em ordem alfabetica e no formato da norma.
\bibliography{referencias}

% Glossario (opcional, 4.2.3.2), em ordem alfabetica
\begin{glossario}
  \termo{Termo}{definição do termo técnico usado no texto.}
  \termo{Outro termo}{definição do outro termo.}
\end{glossario}

% Apendices (opcionais, 4.2.3.3): textos feitos por voce
\apendice{Título do apêndice}
Aqui vai o conteúdo do apêndice, como um questionário ou um roteiro de entrevista.

% Anexos (opcionais, 4.2.3.4): textos de outros autores
\anexo{Título do anexo}
Aqui vai o conteúdo do anexo, como um documento oficial. A referência do anexo fica em nota de rodapé ou numa lista própria.

\end{document}
